\usepackage[brazil]{babel}
\usepackage[T1]{fontenc}
\usepackage[utf8]{inputenc}
\usepackage[hidelinks]{hyperref}
\usepackage[pt-BR]{datetime2}
\usepackage{url}
\usepackage{tikz}
\usetikzlibrary{automata, positioning, arrows, arrows.meta, shapes.geometric, fit, calc}



\usepackage[portuguese,ruled,linesnumbered,algochapter,titlenumbered]{algorithm2e}
\usepackage{listings}
\usepackage{minted}

% Citações (NBR 10520:2023) e referências (NBR 6023:2018), alinhadas à esquerda.
% natbib=true mantém \citet (narrativa) e \citep (entre parênteses);
% \cite é usado no texto sempre de forma narrativa.
% O csquotes vem depois do minted (fvextra), como esse pacote recomenda.
\usepackage{csquotes}
\usepackage[style=abnt,backend=biber,natbib=true,uniquelist=false]{biblatex}
\let\cite\textcite
\addbibresource{references.bib}
% URLs na mesma fonte do texto; permite quebrar URLs longas sem espaçar os caracteres
\urlstyle{same}
\setcounter{biburlnumpenalty}{9000}
\setcounter{biburllcpenalty}{7000}
\setcounter{biburlucpenalty}{8000}

\usepackage{graphicx}
\usepackage{tabularx}
% Compatibilidade: o colortbl (2024+) usa \insert@pcolumn, que não existe no
% pacote array de 2023 instalado no MiKTeX local. Sem efeito em versões novas.
\makeatletter
\providecommand\insert@pcolumn{\insert@column}
\makeatother
% Colunas de tabela alinhadas à esquerda (evita espaços esticados em colunas estreitas).
% O nome P evita conflito com a coluna L já definida em tcc.cls.
\newcolumntype{P}[1]{>{\raggedright\arraybackslash}p{#1}}
\renewcommand{\tabularxcolumn}[1]{>{\raggedright\arraybackslash}p{#1}}
% Ajustes finos de tipografia e folga extra na quebra de linhas
\usepackage[nopatch=footnote]{microtype}
\setlength{\emergencystretch}{2em}
\usepackage{blindtext}
\usepackage{lipsum}
\usepackage{float}
\usepackage{amsthm}
\usepackage{xcolor}
\graphicspath{{./imagens/}}

